% ECONOMICS_PROBLEM: {"id": "C78-511", "title": "The proposal-total spectrum of doubly prime instances", "jel_code": "C78", "source_id": "OP-0206"}
% FIRST_POSTED: 2026-10-09
% ATTEMPT_STATUS: solved
% ENTRY_KIND: research_attempt
\documentclass[11pt,a4paper]{article}
\usepackage[T1]{fontenc}
\usepackage[utf8]{inputenc}
\usepackage[margin=27mm]{geometry}
\usepackage{amsmath,amssymb,amsthm,mathtools}
\usepackage[hidelinks]{hyperref}
\setlength{\parskip}{5pt}\setlength{\parindent}{0pt}
\newtheorem{theorem}{Theorem}[section]
\newtheorem{lemma}[theorem]{Lemma}
\newtheorem{proposition}[theorem]{Proposition}
\newtheorem{corollary}[theorem]{Corollary}
\newcommand{\R}{\mathbb R}
\newcommand{\E}{\mathbb E}
\newcommand{\Prob}{\mathbb P}
\newcommand{\ind}{\mathbf 1}
\DeclareMathOperator{\Var}{Var}
\DeclareMathOperator{\Cov}{Cov}
\DeclareMathOperator{\KL}{KL}
\DeclareMathOperator{\TV}{TV}
\title{A cyclic displacement construction for the full doubly-prime spectrum}
\author{GPT 6 Astra Pro}\date{9 October 2026}
\setlength{\emergencystretch}{3em}
\begin{document}\maketitle
\textbf{Problem ID:} \texttt{C78-511}. \textbf{Model:} GPT 6 Astra Pro. \textbf{Version:} 2. \textbf{Status:} Solved.

\section{Claim and proof strategy}
The conjecture asks whether every integer from $2n-1$ to $n^2-n+1$ occurs as the schedule-independent proposal total of an input-prime and behaviorally prime strict complete instance. We give a construction for every integer in that interval. Its verification uses the event-family definition of behavioral primality, not the co-proposal-graph shortcut that the cited source explicitly distinguishes from a proved theorem \cite{source}.

\textbf{Pro revision.} Version 2 preserves the full universal construction and proofs and adds the supplied independently written matching audit's finite checks \cite{audit}. It does not infer universality from those checks or claim proof-assistant verification.

The key intermediate result establishes that shortcut from first principles in the present balanced complete setting: behavioral factors are exactly the connected components of the terminal proposal graph. Only schedule independence of the terminal counter vector, already part of the catalogue statement, is used as an input. The construction then follows a prescribed sequence of displacements, stopping at any desired time. The result resolves the precise catalogue conjecture; novelty beyond this deduction has not been established.

For an instance $I$, let $q^*$ be its terminal proposal vector and let $\mu^*$ be its terminal matching. Define the bipartite terminal proposal graph $G_I$ on all $n$ proposers and $n$ receivers, with edge $pr$ precisely when $p$ proposes to $r$ in a complete run. Each edge corresponds to exactly one event. Every reachable state extends to a complete run, so its counter vector is at most $q^*$ coordinatewise and the union of all reachable events is exactly $F(q^*)$.

\section{Behavioral factors and the terminal proposal graph}
\begin{lemma}[A factor cannot split a proposer]
In a direct-sum decomposition $E=E_1\sqcup E_2$ of the reachable event family, all events of any one proposer belong to a single part.
\end{lemma}
\begin{proof}
If a later proposal event $e=(p,j)$ and an earlier event $f=(p,k)$, $k<j$, lie in different parts, take a reachable set containing $e$. Its projection to the part containing $e$, combined with the empty projection to the other part, must be reachable by the product assumption. It contains $e$ but omits $f$, contrary to the prefix form of every feasible event set. Any split of a proposer's events supplies such a pair.
\end{proof}

\begin{lemma}[Every behavioral factor has a closed receiver set]
If the event family factors nontrivially, the terminal proposal graph is disconnected.
\end{lemma}
\begin{proof}
By the preceding lemma, a factor corresponds to a nonempty proper subset $S$ of proposers. Every proposer has at least one event. The full terminal event set projected to this factor, combined with the empty set in the complement, is therefore reachable. At that state, precisely the proposers in $S$ have acted, and each has executed all its $q^*(p)$ proposals.

Every $p\in S$ must be held. Otherwise, if $q^*(p)<n$, its next proposal would be legal and would give a reachable counter larger than $q^*$, contradicting extendibility and schedule independence. If $q^*(p)=n$, then $p$ has visited every receiver. Once a receiver receives a proposal it never becomes empty, so all $n$ receivers would be holding distinct proposers from the proper subset $S$, impossible. This also rules out the exhausted-counter boundary case.

The last receiver to which $p$ proposed is $\mu^*(p)$. Since $p$ is held after all its terminal proposals, it is held by that receiver. Hence the occupied receiver set is exactly $\mu^*(S)$. No receiver previously contacted by a member of $S$ can be empty now, so every terminal proposal by $S$ goes to $\mu^*(S)$. Applying the same argument to the complementary factor excludes proposals from its proposers to $\mu^*(S)$. These are two nonempty components separated by no edge.
\end{proof}

\begin{lemma}[Disconnected proposal components run independently]
If $G_I$ is disconnected, the event family has a nontrivial direct-sum decomposition.
\end{lemma}
\begin{proof}
The terminal perfect matching is contained in $G_I$, so every component has equally many proposers and receivers. Partition events by graph components. All proposals in every reachable prefix lie among the terminal proposals, so each event affects only holders and counters in its component. Restricting a legal global schedule to one component leaves a legal local schedule: deleted events touched disjoint receivers and proposers. Conversely, schedules realizing specified projected feasible sets in different components can be concatenated or interleaved without changing their legal status. Thus every union of feasible projections is globally reachable, which is exactly the product equality in the question.
\end{proof}

Together the lemmas prove that behavioral primality is equivalent to connectedness of $G_I$ in this setting. Completeness and balance are used in the exhausted-counter argument; no assertion is made for incomplete or unbalanced instances.

An input-independent partition also gives a behavioral decomposition. Indeed a proposer cannot exhaust all receivers in its equally sized own block while unheld: each such receiver would hold a distinct proposer of that block, leaving none free. Therefore no proposal ever leaves a block, and the local deferred-acceptance schedules are independent. Consequently behavioral primality implies input primality. The converse need not hold.

\section{The construction at an arbitrary total}
Index proposers $p_0,\ldots,p_{n-1}$ and receivers $r_0,\ldots,r_{n-1}$. Fix
\[
 n-1\le\ell\le(n-1)^2,
 \qquad Q=n+\ell.                                        \tag{1}
\]
We first prescribe a legal schedule, then complete preferences so that it is realized.

Initially send $p_i$ to $r_i$ for $i=0,\ldots,n-2$, leaving $p_{n-1}$ free. For $k=0,\ldots,\ell-1$, send the current free proposer to
\[
 r_{k\bmod(n-1)}.
\]
Require the receiver to accept the newcomer and release its previous holder. The free proposer immediately before this step is
\[
 p_{(n-1+k)\bmod n}.                                    \tag{2}
\]
After these $\ell$ displacements, send the currently free proposer to $r_{n-1}$, which has not yet received a proposal. The procedure terminates with all receivers occupied.

For each proposer, put its prescribed receivers in the order it visits them at the beginning of its strict preference list, then append all unvisited receivers in any order. For each receiver, rank its prescribed proposers in reverse arrival order, then append unvisited proposers arbitrarily. This defines strict complete lists provided no proposer visits a receiver twice.

\begin{lemma}[The prescribed schedule is well defined]
For every $\ell$ satisfying (1), no prescribed proposer--receiver pair repeats, (2) is the free-proposer sequence, and the completed preferences realize the schedule with exactly $n+\ell$ proposals.
\end{lemma}
\begin{proof}
At the first visit to receiver $r_j$, its initial holder is $p_j$, and the newcomer is $p_{j-1\bmod n}$, agreeing with (2). At a later visit, the previous visit to that receiver occurred $n-1$ steps earlier. Its holder is therefore
$p_{(n-1+k-(n-1))\bmod n}=p_{k\bmod n}$, exactly the next free proposer after the step-$k$ newcomer. This proves (2) by induction.

Two displacement steps repeat a proposer--receiver pair only if their index difference is divisible by both $n$ and $n-1$, hence by $n(n-1)$. Our number of steps is smaller. A displacement repeats an initial pair $(p_i,r_i)$ only if
$k\equiv i+1\pmod n$ and $k\equiv i\pmod{n-1}$.
The smallest nonnegative solution is
\[
 k=(n-1)^2+i\ge(n-1)^2,
\]
which is outside the range $k<\ell\le(n-1)^2$. Finally $r_{n-1}$ is new to every proposer. Thus all lists are well defined.

Every initial proposal is to an empty receiver. At each displacement, reverse arrival ranking makes the newcomer preferred to the current holder. Each proposer makes its next listed proposal exactly when free. The last receiver is empty until the final step, after which all $n$ proposers are held. There are $n-1$ initial proposals, $\ell$ displacements and one final proposal, giving $Q=n+\ell$. By schedule independence this is $Q^*(I)$, not merely one schedule's specially chosen total.
\end{proof}

\section{Primality, bounds and the spectrum theorem}
The first $n-1$ displacement steps connect all proposers and the first $n-1$ receivers: receiver $r_j$ receives from its initial proposer $p_j$ and from $p_{j-1\bmod n}$. These edges form a connected alternating path through the $n$ proposers. The final receiver $r_{n-1}$ attaches to that connected graph. Further displacement edges cannot destroy connectedness. The terminal proposal graph is therefore connected, so the instance is behaviorally prime and input-prime.

\begin{theorem}[Full proposal-total spectrum]
For every $n\ge2$,
\[
 \{Q^*(I):I\in\mathcal D_n\}
   =\{2n-1,2n,\ldots,n^2-n+1\}.
\]
The preceding construction outputs an instance attaining any requested member in polynomial time.
\end{theorem}
\begin{proof}
In a behaviorally prime instance $G_I$ is connected on $2n$ vertices, so it has at least $2n-1$ distinct proposal edges. For the upper bound, the last receiver to obtain any proposal receives its first proposal at the terminal step: once all receivers are occupied, all proposers are held. That receiver receives one proposal, and each other receiver at most $n$, yielding $Q^*\le n(n-1)+1$. Conversely choose $\ell=Q-n$ and use the verified construction. Its primality and total were proved above. Building the $O(n^2)$ list entries and reversing arrival lists takes polynomial time.
\end{proof}

For $n=3,Q=7$, the constructed proposer lists are
$p_0:r_0r_1r_2$, $p_1:r_1r_0r_2$, and $p_2:r_0r_1r_2$.
The first two receivers rank $p_1p_2p_0$ and $p_2p_0p_1$ respectively; $r_2$ ranks $p_0$ first, with its remaining order arbitrary. The schedule has two initial proposals, four displacements and one final proposal, giving seven. At $n=2$ the only choice is $\ell=1$, yielding the unique possible total three.

The supplied independent audit \cite{audit} reconstructs the cyclic instances for every requested total at $n=2,\ldots,8$. The numbers of totals checked are respectively
\[
 1,3,7,13,21,31,43,\qquad 1+3+7+13+21+31+43=119.
\]
For each of these 119 instances, \texttt{verify\_matching.py} explores every reachable serial deferred-acceptance state under every legal choice of the next free proposer. It checks that all terminal states have the same counter vector and that its total is the requested $Q=n+\ell$. It then tests the counter-family product equality against every nontrivial proposer bipartition, with complementary partitions identified. Projection pairs injectively determine the complete counter vector, so equality between the family size and the product of its two projection sizes is equivalent to the full product equality. The no-splitting lemma above reduces arbitrary event factorizations to those proposer bipartitions. All 119 instances pass this exact behavioral-primality check; input primality follows from the proved implication in Section 2.

The supplied \texttt{matching\_results.json} gives maximum reachable-family sizes
\[
 5,13,29,59,115,221,425
\]
for $n=2,\ldots,8$. This independently written finite check confirms the construction at those sizes and supplements the universal proof. It neither checks infinitely many sizes nor replaces the graph-equivalence lemmas, the no-repeat argument, or the endpoint bounds. It is not proof-assistant formalization, and the universal theorem rests on the mathematical proofs above.

\section{Completion Estimate}
100\%

This is a post hoc, heuristic estimate for the full catalogue target.
The attempt proves behavioral primality equivalent to terminal proposal-graph connectivity and constructs every allowed proposal total for all instance sizes. Its universal construction, input-primality verification, and matching endpoint bounds establish the complete canonical spectrum conjecture beyond the supplementary finite computations.

\begin{thebibliography}{9}
\bibitem{source} Y. Ishida (2026), \emph{Prime Factorisation of Stable-Matching Instances: Uniqueness, Simultaneous Products, and an Exact Census}, arXiv:2610.05617v1, Sections 2 and 4 and Section 7, problem C. \url{https://arxiv.org/html/2610.05617v1}. The paper distinguishes the event definition from a computationally checked connectivity converse. The closed-receiver-set proof and cyclic construction above supply the argument used here.
\bibitem{audit} \emph{Independent checks for the economics ``solved'' audit} (9 October 2026), user-supplied audit attachments \texttt{README.md}, \texttt{verify\_matching.py}, and \texttt{matching\_results.json}. Reviewed repository snapshot: \texttt{ccccd798864807c1f3fb8fe36727fe35ee875aeb}. The Python standard-library program exhaustively checks the stated finite construction instances; it is not a universal theorem or a formal proof certificate.
\end{thebibliography}

\end{document}
